\documentclass[10pt,a4paper]{article}

\usepackage[T1]{fontenc}
\usepackage[utf8]{inputenc}
\usepackage{lmodern}
\usepackage{amsmath,amssymb,amsthm}
\usepackage{booktabs,array,tabularx}
\usepackage{geometry}
\usepackage{xcolor}
\usepackage{enumitem}
\usepackage{microtype}
\usepackage{hyperref}
\usepackage{fancyhdr}

\geometry{margin=1in}
\hypersetup{
  colorlinks=true,
  linkcolor=blue!55!black,
  citecolor=blue!55!black,
  urlcolor=blue!55!black,
  pdftitle={BLS You: A Zero-Query Rogue-Key Forgery Against the 2P-AISO-AS Chosen-Key Interface},
  pdfauthor={Martin Feussner},
  pdfsubject={Cryptanalytic audit report for the Definition 33 interface in IACR ePrint 2026/1394},
  pdfkeywords={digital signatures, BLS, rogue key, public-key aggregation, 2P-AISO-AS, forgery}
}
\setlist{nosep,leftmargin=*}
\setlength{\parindent}{1em}
\setlength{\parskip}{0.25em}
\setlength{\emergencystretch}{2em}
\Urlmuskip=0mu plus 1mu
\pagestyle{fancy}
\fancyhf{}
\fancyhead[L]{BLS You}
\fancyhead[R]{18 October 2026}
\fancyfoot[C]{\thepage}
\renewcommand{\headrulewidth}{0.3pt}

\newtheorem{theorem}{Theorem}
\newcommand{\GOne}{\mathbb{G}_1}
\newcommand{\GTwo}{\mathbb{G}_2}
\newcommand{\GT}{\mathbb{G}_T}
\newcommand{\Zr}{\mathbb{Z}_r}
\newcommand{\KAg}{\mathsf{KAg}}
\newcommand{\Vrfy}{\mathsf{Vrfy}}

\title{\bfseries BLS You:\\
A Zero-Query Rogue-Key Forgery Against the\\
2P-AISO-AS Chosen-Key Interface}
\author{Martin Feussner\\
\small Selmer Center, University of Bergen\\
\small \texttt{martin.feussner@uib.no}}
\date{18 October 2026}

\begin{document}
\maketitle

\noindent\fcolorbox{black!55}{black!4}{\parbox{0.94\textwidth}{
This audit used the OpenAI model Daybreak Blue at the ultra reasoning level.
OpenAI Codex generated and selected this paper's title and deliberately chose
``BLS You'' as BLS/``bless you'' wordplay tied to the scheme name. OpenAI
Codex performed AI-assisted target reconstruction, cryptanalysis,
implementation, experimentation, review coordination, and technical writing.
The source-backed primary reproduction returned PASS, and a fresh-context AI
agent independently reproduced 32/32 attacks and all controls. A separately
tasked AI adversarial reviewer attempted to disprove the result and returned
\texttt{PASS\_NARROW\_SCOPE}, requiring every claim to remain limited to the
direct-addition Definition~33 interface. Independent human verification was
not performed; it is optional and is not a publication requirement.}}

\begin{abstract}
The security framework in IACR ePrint 2026/1394 requires its underlying
two-party signature scheme with aggregatable public keys to satisfy the
chosen-key experiment in Definition~33. That experiment reveals one honestly
generated public key and then accepts an arbitrary second public key from the
adversary before aggregating the pair. The paper's concrete BLS instantiation
aggregates keys by group addition and places no proof-of-possession or
key-knowledge condition on this interface.

We apply the standard BLS rogue-public-key cancellation attack. After seeing
$X_A$, choose a known nonzero scalar $\beta$ and register
$X_E=\beta P-X_A$. The aggregate is $\beta P$, so for any fresh message
$m^\star$ the adversary outputs $S^\star=\beta H(m^\star)$. The ordinary BLS
pairing equation accepts. The attack makes zero signing-oracle queries, uses
full BLS12-381 parameters, and supplies valid nonidentity subgroup elements.

An exact-verifier Rust harness imported the unchanged paper-linked artifact,
used its random key generation, SHA-256 XMD hash-to-$\GTwo$, direct G1 key
addition, and ordinary pairing verifier, and accepted the forgery in 0.03
seconds with 3,072 KiB peak resident memory on one worker. Changed-message,
changed-signature, honest-key-substitution, encoding, subgroup, and
nonidentity controls passed. The result concerns only Definition~33's
chosen-key interface. It is not a claim against Figure~6, Theorem~1, central
honest \textsf{KeyGen}, a proper DKG, the artifact's normal lifecycle, or any
deployment. The method is established prior art, and no technique novelty is
claimed.
\end{abstract}

\section{Introduction and exact claim}

Gomez-Martinez, Tairi, Moreno-Sanchez, and Schneidewind introduce Two-Party
Asymmetric-Input Solitary-Output Adaptor Signatures (2P-AISO-AS) and give a
BLS-based realization in ePrint 2026/1394~\cite{target}. Their preliminary
syntax requires an underlying two-party signature scheme $DS_2$ with
aggregatable public keys and requires it to satisfy two-party existential
unforgeability. Appendix Definition~33 spells out the required experiment.

That experiment has a strong chosen-key order: the challenger generates one
honest key, gives the honest verification key to the adversary, and then asks
the adversary for the other verification key. The two keys are aggregated
before ordinary signature verification. No proof of possession, proof of
secret-key knowledge, distinct-message rule, key-prefixed message, or
key-dependent coefficient appears in the interface.

This is the classic setting for a BLS rogue-key cancellation. We give the
full equation, execute it through the paper-linked artifact's exact ordinary
BLS verifier, and state the lifecycle boundary explicitly. Our claim is:

\begin{quote}
Definition~33's chosen-key $DS_2$ interface is not existentially
unforgeable when instantiated by the paper's direct-addition BLS key
aggregation: a corrupt participant makes a fresh-message forgery with zero
signing queries.
\end{quote}

We do \emph{not} claim a break of the separate Figure~6 lifecycle,
Theorem~1, central honest \textsf{KeyGen}, a proper distributed key-generation
protocol, the artifact's normal application programs, or a deployed system.
The theorem is conditional on its stated underlying primitives; showing that
the direct-addition BLS instantiation fails the Definition~33 assumption does
not refute that conditional theorem. We also do not claim that rogue-key
cancellation is new. The target-specific result is that the paper expressly
requires Definition~33 while its concrete direct-addition BLS interface
admits the standard attack.

\section{Target and evidence identity}

The target is the retained 83-page PDF of \emph{Adaptor Signatures Meet BLS:
Enabling Efficient Blockchain Applications with Unique Adaptor Signatures},
IACR ePrint 2026/1394, received 8 July and approved 12 July 2026. Its SHA-256
digest is
\begin{center}
\small\ttfamily
93c36cf110ab278be767c367c20ae2035446c1bbffc3e1429551d8b42e473e4f.
\end{center}
The exact text reviewed for Definition~33 has SHA-256
\begin{center}
\small\ttfamily
890fab6e15b5276537881eaefa49e8224226dd3933ce7e018e38ba0e560a5001.
\end{center}

The paper links an Anonymous GitHub artifact~\cite{artifact}. The retained official ZIP has
SHA-256
\begin{center}
\small\ttfamily
f1e70afef59a55e3a852d1a9debfefb7f9250f90679a65faa46ecdf0c0a7b441.
\end{center}
It contains 65 files and 178,502 uncompressed bytes. The mirror exposes no
Git repository identity, branch, commit, or tree, so the ZIP digest and the
complete per-file manifest are the authoritative source pin. We make no Git
revision claim. The relevant unchanged files include:

\begin{center}
\begin{tabularx}{\textwidth}{@{}>{\ttfamily\small}l X@{}}
\toprule
File & SHA-256 \\
\midrule
bls-aiso-as/src/bls.rs &
\ttfamily\small f26362cdb3eb14485e62a6fd7085ffb58cdd0597013f52a39b260de7d49de52e \\
bls-aiso-as/src/common.rs &
\ttfamily\small cd7a2d0498f1f62ed8679e70440c9c3b2057f3cd8094ebd5d87d557a1c188344 \\
test-as/src/main.rs &
\ttfamily\small bef2c4fb6a0f6b2de5fbdbf4fdf390a58ea81daa0a4c6bc4fbd0ef8b67b50d12 \\
test-as-optimized/src/main.rs &
\ttfamily\small e33790341c5ad68095dcb6808b7a0232938d49edf3104a81034faeb57889dadf \\
\bottomrule
\end{tabularx}
\end{center}

The artifact uses the BLS12-381 implementation distributed as
\texttt{irmaseal-curve} 0.1.4. Public keys are G1 points, signatures are G2
points, and messages are hashed to G2 using SHA-256 XMD with domain string
\texttt{BLS12-381-G2:SHA-256}. The ordinary verifier in
\texttt{bls.rs}, lines 110--114, is the endpoint exercised below.

\section{Definition 33 and the BLS endpoint}

Let $\GOne$, $\GTwo$, and $\GT$ have prime order $r$, let $P$ generate
$\GOne$, and let
\[
 e:\GOne\times\GTwo\longrightarrow\GT
\]
be a nondegenerate bilinear pairing. Let $H:\{0,1\}^*\to\GTwo$ be the
artifact's hash-to-curve function. In the artifact's orientation, a secret
$x\in\Zr$ has public key $X=xP$, and a signature on $m$ is
$S=xH(m)$. Ordinary verification accepts exactly when
\begin{equation}
  e(P,S)=e(X,H(m)).                                    \label{eq:verify}
\end{equation}

The concrete two-key aggregation used by the artifact test programs is direct
G1 addition:
\begin{equation}
  \KAg(X_A,X_E)=X_A+X_E.                               \label{eq:kag}
\end{equation}

Definition~33 first initializes an empty query set $Q$, generates the honest
keypair, and reveals $X_A$ to the adversary. Only then does it accept the
adversarial key $X_E$. After an optional signing-oracle phase, it computes
$\bar X=\KAg(X_A,X_E)$ and accepts an adversarial pair $(m^\star,S^\star)$
when $m^\star\notin Q$ and Equation~\eqref{eq:verify} holds under $\bar X$.
The attack below never invokes the signing oracle, so $Q$ remains empty.

\section{Zero-query rogue-key forgery}

\begin{theorem}
The direct-addition BLS instantiation fails Definition~33 with probability
one, apart from avoidable choices of an identity key, using zero signing
queries and constant group work.
\end{theorem}

\begin{proof}
After receiving the honest public key $X_A$, choose any known
$\beta\in\Zr\setminus\{0\}$ for which $\beta P\neq X_A$, and return
\begin{equation}
  X_E=\beta P-X_A.                                     \label{eq:rogue}
\end{equation}
The condition is publicly testable. If one candidate produces the identity,
select another nonzero $\beta$. Because $X_A$ and $\beta P$ are elements of
the prime-order subgroup and that subgroup is closed under subtraction,
$X_E$ is a valid subgroup element. The explicit inequality makes $X_E$
nonidentity. Equation~\eqref{eq:kag} now gives
\begin{equation}
  \bar X=X_A+X_E=\beta P,                              \label{eq:cancel}
\end{equation}
which is also nonidentity because $\beta\neq0$.

Choose any fresh message $m^\star$ and output
\begin{equation}
  S^\star=\beta H(m^\star).                            \label{eq:forge}
\end{equation}
The adversary makes no signing query. Bilinearity and
Equation~\eqref{eq:cancel} give
\[
 e(P,S^\star)
 =e(P,\beta H(m^\star))
 =e(\beta P,H(m^\star))
 =e(\bar X,H(m^\star)),
\]
so the ordinary verifier accepts. The adversary knows the aggregate secret
$\beta$ but need not know the discrete logarithm $\beta-x_A$ of $X_E$.
Definition~33 asks only for the public key and contains no knowledge check.
\end{proof}

The work is one scalar multiplication and one subtraction in G1 to construct
the corrupt key, one hash to G2, and one scalar multiplication in G2 to forge.
Memory is constant. There is no search, reduced parameter, collision,
preimage, malformed encoding, parser behavior, weak honest key, fault,
side-channel, or implementation bug.

\section{Source-backed execution and controls}

The primary Rust harness imports the unchanged \texttt{bls-aiso-as} crate. It
calls \texttt{BLSKeyPair::new()} for the honest key, selects $\beta=42$ (or
$43$ in the publicly detected identity edge case), constructs
Equation~\eqref{eq:rogue}, applies direct G1 addition, calls the artifact's
\texttt{hash\_to\_g2}, and sends Equation~\eqref{eq:forge} to the public
\texttt{bls::verify}. It contains no call to a signing function.

The full BLS12-381 execution returned PASS. The direct attack run took 0.03
seconds and 3,072 KiB maximum resident memory; a sequential repeat took 0.04
seconds with the same memory and byte-identical standard output. Both used one
worker and a 1 GiB virtual-memory limit. A clean locked offline build-and-run
took 1.35 seconds and 120,320 KiB. The source output recorded
\texttt{kag\_equals\_beta\_g1=true}, zero signing queries, a fresh message,
and acceptance by the ordinary artifact verifier.

The harness also establishes:
\begin{itemize}
  \item the honest, corrupt, and aggregate G1 keys are nonidentity and
        torsion-free;
  \item the forged signature is computed in the G2 prime-order subgroup;
  \item changing the message rejects;
  \item adding the G2 generator to the signature rejects;
  \item replacing the honest key rejects; and
  \item malformed all-\texttt{ff} compressed G1 and G2 strings are rejected
        by the concrete checked decoding boundary.
\end{itemize}

A separately written fresh-context harness ran 32 randomized honest-key
trials through the same unchanged verifier and returned 32/32 attack
successes with all changed-message, changed-signature, honest-key-substitution,
nonidentity, subgroup, checked-encoding, and malformed-encoding controls
passing. Its cryptographic checks took 1.16381577 seconds; the clean locked
compile-and-run took 102.98 seconds wall and 400,048 KiB maximum resident
memory with one worker. The sealed \texttt{REVIEW\_ARTIFACTS.sha256} file has
SHA-256
\texttt{b9c45ff8afea908c2732690215ad80f5acb0e67b4c49567cb304c449ff798a8c};
the state handoff has SHA-256
\texttt{b089529d0146daeb37bd50f370411c311012283a25d1ca090aaaf128eb67d22e}.

The sealed post-evidence adversarial review returned
\texttt{PASS\_NARROW\_SCOPE}. It independently checked the key-choice order,
validity and ownership distinction, aggregation semantics, unchanged verifier,
freshness, controls, practical resources, profile pin, prior-art attribution,
and lifecycle boundary. Its manifest has SHA-256
\texttt{f4558bb0b73c6e0a6ae7c1b8b400ebdc2f3cabde0ec71f93cbc7c01ba54eb065},
and its handoff has SHA-256
\texttt{b7218c2be45c0683a1243a639e7be3d510db30018efa230ccaefb6d226b6a2db}.

\section{Lifecycle boundary}

The relevant interfaces must remain separate:

\begin{center}
\begin{tabularx}{\textwidth}{@{}p{0.24\textwidth}p{0.30\textwidth}X@{}}
\toprule
Interface & Key provenance & Claim \\
\midrule
Definition 33 & Honest $X_A$ is revealed before the adversary chooses
$X_E$; no proof of possession & Broken by the zero-query forgery \\
Figure 6 & Both participant keypairs come from central honest
\textsf{KeyGen} & Excluded from this claim \\
Theorem 1 & Conditional construction theorem whose assumptions include the
required security of the underlying $DS_2$ & Not refuted by this result \\
Artifact programs & Both keypairs are independently generated with
\texttt{BLSKeyPair::new()} & Excluded from this claim \\
Proper DKG or authenticated registration & Requires a separately specified
knowledge and binding guarantee & Excluded and not evaluated \\
Deployments & Deployment-specific registration and lifecycle rules & No
deployment claim \\
\bottomrule
\end{tabularx}
\end{center}

Figure~6 and the normal artifact programs do not let the corrupt participant
replace its key after seeing the honest key. Footnote~7 describes central key
generation as a syntactic simplification that could be realized by DKG. That
is materially different from Definition~33. The source execution is therefore
an exact-verifier diagnostic for the required chosen-key interface, not an
end-to-end artifact-lifecycle demonstration.

This boundary does not make the Definition~33 result merely an implementation
observation. The interface is printed as the required security notion for the
underlying $DS_2$, and the cancellation follows from its key order and the
concrete aggregation equation. It does mean that every title, abstract, table,
and summary must name Definition~33 rather than saying without qualification
that all of 2P-AISO-AS or a deployment is broken.

\section{Prior art and history search}

The method is established. Boneh, Gentry, Lynn, and Shacham discuss BLS
aggregate signatures and the chosen-public-key problem, using distinct
messages or public-key-prefixed inputs as a defense~\cite{bgls}. Bellare and
Neven formalize multisignatures in the plain public-key model, where malicious
keys may depend on honest keys~\cite{bn}. Boneh, Drijvers, and Neven give the
same cancellation explicitly: choose a key equal to a known generator power
times the inverse of the honest key, then sign under the known aggregate
exponent~\cite{bdn}. The current CFRG BLS draft describes Basic, Augmentation,
and Proof-of-Possession variants and warns that fast common-message aggregate
verification requires a valid proof for every public key~\cite{cfrg}.

Against the Definition~33 chosen-key $DS_2$ interface of ePrint 2026/1394,
we apply the standard BLS rogue-public-key cancellation attack and obtain a
zero-query fresh-message forgery at full BLS12-381 parameters. We located no
earlier public report applying this attack to that exact interface in searches
completed through 18 October 2026. The attack technique is prior art; no
technique novelty is claimed. This is a dated, bounded search result and does
not exclude private, deleted, embargoed, unindexed, or later disclosures.

The final history recheck returned \texttt{PASS\_FINAL\_HISTORY}. Its bundle
manifest has SHA-256
\texttt{796daa3485222c38b14b3e6c6b2e6c7a4514df09efb60c8bbb31b67ada0e25f3};
the 74-source manifest has SHA-256
\texttt{a12410ab6c67aee1d3a2fbd9caf2acc2667007ccfaa4d1512f23ebfd061f9e95};
and the handoff has SHA-256
\texttt{847ad6ca5617a2491728e178b957cdc2a096cf6beb9ba0a9c6c7919f072481e3}.

\section{Repair directions and limitations}

Public-key subgroup and nonidentity checks are necessary but do not stop this
attack: $X_E$ already passes them. Standard design directions include:
\begin{itemize}
  \item require and verify a proof of possession or proof of secret-key
        knowledge before accepting either contribution;
  \item augment each signed message with the corresponding public key;
  \item use list-bound, key-dependent aggregation coefficients as in the BDN
        line of work; or
  \item replace the chosen-key interface with authenticated honest generation
        or a DKG that proves possession and binds every contribution.
\end{itemize}
Each option changes an interface or equation used by the paper. Its fit with
the two-party protocol, adaptor operations, uniqueness properties, and proof
must be analyzed afresh. This audit has not implemented or validated a repair.

The evidence has the following limits:
\begin{itemize}
  \item the claim is exactly the Definition~33 chosen-key $DS_2$ interface;
  \item Theorem~1, Figure~6, and the normal artifact key lifecycles are outside
        the attack claim;
  \item the artifact exposes typed group points rather than a public
        adversarial byte-registration API;
  \item the anonymous mirror provides a hash-pinned ZIP but no Git commit;
  \item the history statement is bounded to the public sources searched
        through 18 October 2026; and
  \item all required technical, adversarial, independent-reproduction,
        final-history, and publication checks were completed before
        classification.
\end{itemize}

\section{Conclusion}

Definition~33 reveals an honest BLS key before accepting an unproved corrupt
key and then aggregates by direct addition. Standard rogue-key cancellation
therefore makes the aggregate key's discrete logarithm known to the adversary.
A fresh BLS signature follows with zero signing queries, and the unchanged
full-parameter artifact verifier accepts it. This result is precise but
narrow: it breaks the printed chosen-key interface, while Theorem~1 and every
honest-key or proper-DKG lifecycle remain outside the attack claim.

\section*{Acknowledgements}

The computations were performed on the Norwegian Research and Education Cloud (NREC), using resources provided by the University of Bergen and the University of Oslo.

\begin{thebibliography}{9}

\bibitem{target}
J. Gomez-Martinez, E. Tairi, P. Moreno-Sanchez, and C. Schneidewind.
\newblock Adaptor Signatures Meet BLS: Enabling Efficient Blockchain
Applications with Unique Adaptor Signatures.
\newblock IACR Cryptology ePrint Archive, Report 2026/1394, 2026.
\newblock \url{https://eprint.iacr.org/2026/1394}.

\bibitem{artifact}
Gomez-Martinez et al.
\newblock Paper-linked 2P-AISO-AS artifact, source snapshot dated 7 May 2026.
\newblock \url{https://anonymous.4open.science/r/Adaptor_Signatures_Meet_BLS-ECD0/}.

\bibitem{bgls}
D. Boneh, C. Gentry, B. Lynn, and H. Shacham.
\newblock Aggregate and Verifiably Encrypted Signatures from Bilinear Maps.
\newblock In \emph{EUROCRYPT 2003}, LNCS 2656, pages 416--432, 2003.
\newblock \url{https://eprint.iacr.org/2002/175}.

\bibitem{bn}
M. Bellare and G. Neven.
\newblock Multi-signatures in the Plain Public-Key Model and a General Forking
Lemma.
\newblock In \emph{ACM CCS 2006}, pages 390--399, 2006.

\bibitem{bdn}
D. Boneh, M. Drijvers, and G. Neven.
\newblock BLS Multi-Signatures With Public-Key Aggregation.
\newblock 2018.
\newblock \url{https://crypto.stanford.edu/~dabo/pubs/papers/BLSmultisig.html}.
\newblock See also IACR ePrint 2018/483,
\emph{Compact Multi-Signatures for Smaller Blockchains}.

\bibitem{cfrg}
D. Boneh et al.
\newblock BLS Signatures, Internet-Draft
\texttt{draft-irtf-cfrg-bls-signature-07}, 6 July 2026.
\newblock \url{https://datatracker.ietf.org/doc/draft-irtf-cfrg-bls-signature/}.

\end{thebibliography}

\end{document}
